\chapter{Spasticity}

\section*{Definition}
Spasticity is a velocity-dependent increase in muscle tone due to hyperexcitability of the stretch reflex, accompanied by exaggerated deep tendon reflexes and clonus. It is a component of upper motor neuron (UMN) syndrome and is distinguished from rigidity by its velocity dependence and association with weakness and Babinski sign.

\section*{Pathophysiology}
Spasticity results from loss of descending inhibitory control from the corticospinal and reticulospinal tracts onto spinal alpha motor neurons and interneurons. Reduced GABAergic and glycinergic inhibition at the spinal level lowers the threshold for the monosynaptic stretch reflex arc. Excessive excitatory input from the reticulospinal tract (which becomes hyperactive after corticospinal damage) further increases motor neuron excitability. The velocity-dependent nature arises because faster stretch produces more rapid Ia afferent firing, overwhelming the impaired inhibitory mechanisms.

\section*{Causes}
\begin{itemize}
  \item \textbf{Stroke:} Most common cause of spasticity; upper limb flexors and lower limb extensors affected
  \item \textbf{Multiple sclerosis:} Spastic paraparesis with flexor spasms
  \item \textbf{Spinal cord injury:} Below the level of injury; complete or incomplete lesions
  \item \textbf{Cerebral palsy:} Spastic diplegia, hemiplegia, or quadriplegia from perinatal injury
  \item \textbf{Traumatic brain injury:} Diffuse axonal injury or focal contusions
  \item \textbf{Amyotrophic lateral sclerosis:} Combined UMN and LMN signs
  \item \textbf{Hereditary spastic paraparesis:} Genetic degenerative disorder of the corticospinal tract
  \item \textbf{Adrenoleukodystrophy, adrenomyeloneuropathy:} Metabolic causes of spastic paraparesis
  \item \textbf{Structural:} Cervical myelopathy, spinal cord tumor, syringomyelia, transverse myelitis
\end{itemize}

\section*{When to See a Doctor}
See your doctor if:
\begin{itemize}
  \item New-onset spasticity after stroke, spinal cord injury, or neurological diagnosis
  \item Spasticity causing functional limitation (impaired ambulation, difficulty with hygiene)
  \item Painful muscle spasms interfering with sleep or daily activities
  \item Development of contractures or pressure ulcers from muscle tightness
\end{itemize}

\section*{Related Symptoms}
\begin{itemize}
  \item Velocity-dependent resistance to passive stretch
  \item Clasp-knife phenomenon: initial resistance that suddenly gives way
  \item Clonus: rhythmic, oscillatory contraction after sustained stretch (ankle > patella)
  \item Hyperreflexia and extensor plantar response (Babinski sign)
  \item Flexor or extensor spasms
  \item Weakness, loss of dexterity, and fatigue
\end{itemize}
\textbf{Important:} Spasticity can be both beneficial (assists with weight-bearing in weak legs) and harmful (impairs function, causes pain and contractures); treatment goals must be individualized.

\section*{Self-Care / Home Management}
\begin{itemize}
  \item Daily passive range-of-motion exercises to prevent contractures
  \item Slow, sustained stretching of spastic muscles (hold for 30--60 seconds)
  \item Proper positioning in seated and lying positions to reduce spastic triggers
  \item Avoidance of noxious stimuli (urinary tract infection, pressure ulcers, ingrown toenails) that exacerbate spasticity
  \item Loose-fitting clothing over spastic limbs
\end{itemize}

\section*{How Your Doctor Will Evaluate You}
\begin{itemize}
  \item Clinical assessment: Ashworth Scale or Modified Ashworth Scale (MAS) for spasticity grading
  \item Tardieu Scale: velocity-dependent assessment with measurement at slow and fast speeds
  \item Examination for clonus, hyperreflexia, and Babinski sign
  \item Identification and elimination of spasticity triggers (pain, infection, constipation)
  \item MRI of brain or spine to identify the underlying cause if not already known
  \item Instrumented gait analysis for objective measurement in research settings
\end{itemize}

\section*{Treatment Options}
\begin{itemize}
  \item Physical and occupational therapy: stretching, constraint-induced movement therapy, neurodevelopmental techniques
  \item Oral medications: Baclofen (GABA-B agonist), tizanidine (alpha-2 agonist), gabapentin, dantrolene, diazepam
  \item Focal therapy: Botulinum toxin A injections for focal spasticity (upper limb flexors, ankle plantar flexors)
  \item Intrathecal baclofen pump for generalized or severe spasticity refractory to oral therapy
  \item Phenol nerve blocks for severe focal spasticity (axillary, musculocutaneous, obturator nerves)
  \item Surgical management: Selective dorsal rhizotomy (cerebral palsy), tendon lengthening, tendon transfer, neurotomy
  \item Orthoses for positioning and contracture prevention (e.g., ankle-foot orthosis, hand splints)
\end{itemize}

\clearpage
